\subsection{Exact rational-input classification without commensurate energies}
\label{sec:rationallogendpoint}
Within the full triplet-failure family, the endpoint criterion admits a
terminating exact decision procedure even when the energy gaps are
incommensurate. The additional ingredient is Baker's linear-independence
theorem for logarithms, rather than a polynomial substitution for the whole
Gibbs curve. This is a consequence of an established transcendence theorem;
no new number-theoretic independence claim is made.

\begin{theorem}[Rational-input endpoint decision]
\label{thm:rationallogendpoint}
Consider a positive reversible four-state generator with rational entries,
distinct ordered Gibbs bath weights and positive finite bath temperature.
Suppose all original triplet predicates fail. Ordinary asymptotic cooling
in relative entropy or weighted Euclidean distance can be decided by a
terminating exact procedure using rational arithmetic,
integer factorization, real algebraic root isolation and rational logarithm
enclosures. Commensurate energy gaps are not required. A repeated slow
eigenvalue is handled by the preceding degenerate-boundary classification,
which excludes an effect in these distances. When the slow mode is simple,
the same decision applies asymptotically to every fixed norm.
The procedure concerns this full-triplet family;
it is not a classifier for arbitrary unconstrained rational generators or
for transient distance crossings.
\end{theorem}
\begin{proof}
The unique normalized stationary law $q$ is rational, since it solves a
rational linear system. Equivalently one may supply rational ordered bath
weights and rational downward rates, construct the reverse rates by detailed
balance, and check the hypotheses exactly. In bath energy units take
$E_i=-\log q_i$; an additive energy shift does not affect the criterion.
All logarithms below are the real logarithms of positive numbers.

Let $r$ be the largest root of the rational nonzero characteristic cubic
$p(x)=\det(xI-M)/x$. Determine its multiplicity exactly. If it is repeated,
the full-triplet degeneracy theorem returns absence of an effect. Otherwise
write
\[
\tau=\operatorname{tr}M,\qquad
\sigma=\frac{\tau^2-\operatorname{tr}(M^2)}2,\qquad
K_r=M^2-(\tau-r)M+(\sigma-\tau r+r^2)I.
\]
Define the rational vectors and matrix
\[
a=\tfrac14\mathbf1-q,\qquad h=D^{-1}a,\qquad
T=I-\tfrac14\mathbf1\mathbf1^\top,
\quad D=\operatorname{diag}(q),
\]
and the row vector
\[
c^\top=\tfrac14 h^\top K_r T\in\mathbb Q(r)^{1\times4}.
\]
The endpoint scalar from the spectral-projector test is
\begin{equation}\label{eq:rationalprimelog}
\mathcal R(r)=\sum_{i=1}^4c_i\log q_i.
\end{equation}
Indeed $\pi'(0)=\tfrac14T(\log q_i)_i$, and on the zero-sum probability
subspace the slow spectral projector is $K_r/p'(r)$. Since $\pi'(0)$
has zero sum, this suffices even though the polynomial need not annihilate
the stationary direction on the full space. For an oriented slow left mode
$w$, the projector identity gives
\[
\frac{\mathcal R(r)}{p'(r)}
=\frac{(w^\top a)(w^\top\pi'(0))}{w^\top Dw}.
\]
Here $p'(r)>0$ for a simple largest nonzero root, including when the two
faster roots coincide, and $w^\top a=b_2(0)<0$ by the contrast theorem.
Thus the single-endpoint criterion says cooling occurs exactly when
$\mathcal R(r)>0$; equality means absence, not an undecided threshold.

Factor the numerators and denominators of the four positive rationals $q_i$.
Let $\mathcal P$ be the resulting finite set of primes and write
\[
q_i=\prod_{\ell\in\mathcal P}\ell^{v_{\ell i}},
\qquad v_{\ell i}\in\mathbb Z.
\]
Then
\[
\mathcal R(r)=\sum_{\ell\in\mathcal P}\gamma_\ell\log\ell,
\qquad
\gamma_\ell=\sum_i v_{\ell i}c_i\in\mathbb Q(r).
\]
Each $\gamma_\ell$ has an exact zero test in the real algebraic field
specified by the isolated root $r$.

The precise external theorem needed is the following form of Baker's result:
if logarithms of nonzero algebraic numbers, together with $2\pi i$, are
linearly independent over $\mathbb Q$, those logarithms are linearly
independent over $\overline{\mathbb Q}$.
This is stated in the opening paragraph of
\href{https://doi.org/10.1112/S0025579300008068}{A.~Baker,
\emph{Linear forms in the logarithms of algebraic numbers (II)},
Mathematika 14 (1967), 102--107}, referring to the original
\href{https://doi.org/10.1112/S0025579300003971}{1966 paper,
Mathematika 13, 204--216}.
For distinct positive primes, rational independence of their real logarithms
follows from unique prime factorization: clear the denominators of a rational
relation and exponentiate. Adjoining $2\pi i$ preserves that independence
because all the prime logarithms are real. Baker's theorem therefore gives
\[
\mathcal R(r)=0\quad\Longleftrightarrow\quad
\gamma_\ell=0\ \text{for every }\ell\in\mathcal P.
\]
This is linear independence over algebraic coefficients, not algebraic
independence of the logarithms. There is no branch ambiguity: only real
logs of primes greater than one are used. A rational factor equal to one
has empty factorization and contributes zero, so $\log1$ is never included
in the independent set.

If every coefficient vanishes, return equality and absence. Otherwise the
sum is a nonzero real number. Refine the rational isolating interval for
$r$, evaluate the algebraic coefficients by directed rational interval
arithmetic, and enclose each logarithm by a convergent rational series.
For example, for a prime $\ell$ put $z=(\ell-1)/(\ell+1)\in(0,1)$.
The partial sum and a valid tail bound are
\[
L_N=2\sum_{j=0}^{N}\frac{z^{2j+1}}{2j+1},\qquad
0<\log\ell-L_N\leq
\frac{2z^{2N+3}}{(2N+3)(1-z^2)}.
\]
All quantities on the right are rational and the bound tends to zero.
Refining all intervals produces an enclosure of $\mathcal R(r)$ whose
width tends to zero. Its nonvanishing guarantees that eventually the
enclosure excludes zero, deciding its sign. No unproved transcendental
equality oracle or floating-point tolerance enters the decision.
\end{proof}

The theorem establishes termination, not a practical complexity bound.
Integer factorization and very small nonzero logarithmic forms can be
expensive; no uniform precision or runtime promise is made. It also does
not remove the full-generator spectral dependence: isolating the one slow
cubic root remains necessary, although a full slow eigenvector and a scan
of initial temperatures are avoided. For arbitrary nonrational input data,
the analytic scalar criterion remains valid but this rational-input
termination argument does not apply.

The implementation is \path{research/exact_rational_triplet_classifier.py}.
It checks the full premise exactly, factors the supplied rational weights,
and uses the selected real algebraic embedding throughout. Range reduction
to logarithms of numbers between one and two accelerates the rational
series. An optional resource limit returns an unresolved result, never
absence. With no such limit the theorem gives termination. The original
counterexample, whose energy gaps are not all commensurate, is classified
as cooling by an exact positive endpoint enclosure. An independent
implementation audit includes commensurate comparisons with the Sturm
classifier, scale invariance, degenerate cases and invalid-input rejection.
